\documentclass[pdflatex, 11pt]{beamer}

\usepackage[utf8]{inputenc}
\usepackage[T1]{fontenc}
\usepackage[italian]{babel}
\usepackage{calligra}
\usepackage{graphicx}
\usepackage{tikz}
\usepackage{epstopdf}
\usepackage{spot}
\usepackage{multirow}
\usepackage{algorithm}% http://ctan.org/pkg/algorithms
\usepackage{algpseudocode}% http://ctan.org/pkg/algorithmicx

\graphicspath{{img/}}

\definecolor{links}{HTML}{2A1B81}
\hypersetup{colorlinks,linkcolor=links,urlcolor=links}

\mode<presentation>{
  \usetheme{Boadilla}
  \usecolortheme{beaver}
  \setbeamercovered{transparent=50}
  \setbeamertemplate{navigation symbols}{}
  \setbeamertemplate{itemize item}{\checkmark}
  \setbeamertemplate{itemize subitem}{-}
  \setbeamertemplate{enumerate items}[default]
}

\usebackgroundtemplate{
  \begin{tikzpicture}
    \node[opacity=0.05] {\includegraphics[scale=1]{img/logoUnifi.eps}};
  \end{tikzpicture}
}

\setspotlightstyle{rectangle}

\newcommand{\var}[1]{{\ttfamily#1}}% variable

\newenvironment<>{varblock}[2][\textwidth]{%
	\setlength{\textwidth}{#1}
	\begin{actionenv}#3%
		\def\insertblocktitle{#2}%
		\par%
		\usebeamertemplate{block begin}}
	{\par%
		\usebeamertemplate{block end}%
	\end{actionenv}}


\begin{document}

\title[Motif Finding]{Analisi di algoritmi per il Motif Finding}
\date[11 Dicembre 2015]{\flushright 11 Dicembre 2015}
\institute[Università di Firenze]{
%  \includegraphics[width=2cm]{img/logoUnifi.eps}\\
%  \includegraphics[width=2cm]{img/logoUnifi.png}\\
%  \includegraphics[]{img/logoUnifiSmall.png}\\
  \includegraphics[width=5cm]{img/logoUnifiName.eps}
}

\author[Papini - Bani]{
  \begin{center}
    \begin{tabular}{lr}
      Tommaso \textsc{Papini}&Gabriele \textsc{Bani}\\
      \href{mailto:tommaso.papini1@stud.unifi.it}{tommaso.papini1@stud.unifi.it}&
      \href{mailto:gabriele.bani@stud.unifi.it}{gabriele.bani@stud.unifi.it}
    \end{tabular}
  \end{center}
}

\titlegraphic{
  \vspace{-0.5cm}
  \tiny
  \href{http://creativecommons.org/licenses/by-sa/4.0/}{\includegraphics[width=1cm]{img/logoCC.png}}
  This work is licensed under a
  \href{http://creativecommons.org/licenses/by-sa/4.0/}{Creative
    Commons Attribution-ShareAlike 4.0 International License}.
}

\begin{frame}[plain]
	\titlepage
\end{frame}

\begin{frame}
	\frametitle{Un po' di background}
	DNA:
	\begin{itemize}
		\item sequenza di nucleotidi
		\item 4 tipi di nucleotide: A, T, C, G
		\item \textit{l-mer}: sottosequenza di DNA di lunghezza $l$
	\end{itemize}
	\begin{block}{Motifs}
		In biologia può essere necessario ricavare certe sequenze di DNA ``nascoste''
		\begin{itemize}
			\item pattern di nucleotidi ripetuti (l-mer)
			\item utili a capire determinati comportamenti biologici
			\begin{itemize}
				\item sequenze di attivazione di geni specifici
			\end{itemize}
		\end{itemize}
	\end{block}
\end{frame}

\begin{frame}
	\frametitle{Il problema del Motif Finding}
	Il problema del Motif Finding consiste nel ricavare un set di $t$ \textbf{l-mer} da un insieme di $t$ sequenze di DNA.
	\begin{block}{Input}
		\begin{itemize}
			\item $DNA$: matrice di nucleotidi $t \times n$
			\begin{itemize}
				\item $t$ sequenze di DNA
				\item ognuna di lunghezza $n$
			\end{itemize}
			\item $l$: lunghezza del motif cercato
		\end{itemize}
	\end{block}
	\begin{block}{Output}
		\begin{itemize}
			\item $s=(s_1,s_2,\dots,s_t)$: lista di $t$ posizioni iniziali di l-mer il più simili tra loro
		\end{itemize}
	\end{block}
\end{frame}

\begin{frame}{Un primo esempio}
	\begin{center}
		CGGGGCT\spot<2->[fill=yellow]{ATGCAACT}GGGTCGTCACATTCCCCTTTCGATA\\
		TTTGAGGGTGCCCAATAA\spot<2->[fill=yellow]{ATGCAACT}CCAAAGCGGACAAA\\
		GG\spot<2->[fill=yellow]{ATGCAACT}GATGCCGTTTGACGACCTAAATCAACGGCC\\
		AAGG\spot<2->[fill=yellow]{ATGCAACT}CCAGGAGCGCCTTTGCTGGTTCTACCTG\\
		AATTTTCTAAAAAGATTATAATGTCGGTCC\spot<2->[fill=yellow]{ATGCAACT}TC\\
		CTGCTGTACAACTGAGATCATGCTGC\spot<2->[fill=yellow]{ATGCAACT}TTCAAC\\
		TACATGATCTTTTG\spot<2->[fill=yellow]{ATGCAACT}TGGATGAGGGAATGATGC
	\end{center}
\end{frame}

\begin{frame}{Mutazioni random}
	\begin{center}
		CGGGGCT\spot<1->[fill=yellow]{ATcCAgCT}GGGTCGTCACATTCCCCTTTCGATA\\
		TTTGAGGGTGCCCAATAA\spot<1->[fill=yellow]{ggGCAACT}CCAAAGCGGACAAA\\
		GG\spot<1->[fill=yellow]{ATGgAtCT}GATGCCGTTTGACGACCTAAATCAACGGCC\\
		AAGG\spot<1->[fill=yellow]{AaGCAACc}CCAGGAGCGCCTTTGCTGGTTCTACCTG\\
		AATTTTCTAAAAAGATTATAATGTCGGTCC\spot<1->[fill=yellow]{tTGgAACT}TC\\
		CTGCTGTACAACTGAGATCATGCTGC\spot<1->[fill=yellow]{ATGCcAtT}TTCAAC\\
		TACATGATCTTTTG\spot<1->[fill=yellow]{ATGgcACT}TGGATGAGGGAATGATGC
	\end{center}
	Come trovare l'l-mer più simile tra tutti?
\end{frame}

\begin{frame}{Allineamento}
	\tiny\begin{align*}
		CGGGGCT&\spot<1->[fill=yellow]{ATcCAgCT}GGGTCGTCACATTCCCCTTTCGATA\\
		TTTGAGGGTGCCCAATAA&\spot<1->[fill=yellow]{ggGCAACT}CCAAAGCGGACAAA\\
		GG&\spot<1->[fill=yellow]{ATGgAtCT}GATGCCGTTTGACGACCTAAATCAACGGCC\\
		AAGG&\spot<1->[fill=yellow]{AaGCAACc}CCAGGAGCGCCTTTGCTGGTTCTACCTG\\
		AATTTTCTAAAAAGATTATAATGTCGGTCC&\spot<1->[fill=yellow]{tTGgAACT}TC\\
		CTGCTGTACAACTGAGATCATGCTGC&\spot<1->[fill=yellow]{ATGCcAtT}TTCAAC\\
		TACATGATCTTTTG&\spot<1->[fill=yellow]{ATGgcACT}TGGATGAGGGAATGATGC
	\end{align*}
\end{frame}

\begin{frame}{Profilo e Consenso}
	\begin{center}
		\begin{tabular}{l l l l l l l l l l}
			\multirow{8}{*}{\textbf{Allineamento}} & & A & T & C & C & A & G & C & T\\
			& & G & G & G & C & A & A & C & T\\
			& & A & T & G & G & A & T & C & T\\
			& & A & A & G & C & A & A & C & C\\
			& & T & T & G & G & A & A & C & T\\
			& & A & T & G & C & C & A & T & T\\
			& & A & T & G & G & C & A & C & T\\
			\hline
			\multirow{4}{*}{\textbf{Profilo}} & \textbf{A} & 5 & 1 & 0 & 0 & 5 & 5 & 0 & 0\\
			& \textbf{T} & 1 & 5 & 0 & 0 & 0 & 1 & 1 & 6\\
			& \textbf{G} & 1 & 1 & 6 & 3 & 0 & 1 & 0 & 0\\
			& \textbf{C} & 0 & 0 & 1 & 4 & 2 & 0 & 6 & 1\\
			\hline
			\textbf{Consenso} & & A & T & G & C & A & A & C & T 
		\end{tabular}
	\end{center}
\end{frame}

\begin{frame}{Score}
	Come definire la ``bontà'' di un set di l-mer?
	\begin{block}{Funzione score}
		Si definisce una funzione score sul vettore $s=(s_1,s_2,\dots,s_t)$ di posizioni iniziali:
		\begin{equation*}
			Score(s,DNA)=\sum_{j=1}^{l}M_{P(s)}(j)
		\end{equation*}
		dove
		\begin{itemize}
			\item $P(s)$: matrice profilo su $s$
			\item $M_{P(s)}(j)$: elemento massimo nella colonna $j$-esima di $P(s)$
		\end{itemize}
	\end{block}
	Si cerca il set di posizioni iniziali $s$ che massimizzi $Score(s,DNA)$!
\end{frame}

\begin{frame}{Score: l'esempio di prima}
	\begin{center}
		\begin{tabular}{l l l l l l l l l l}
			\multirow{8}{*}{\textbf{Allineamento}} & & A & T & C & C & A & G & C & T\\
			& & G & G & G & C & A & A & C & T\\
			& & A & T & G & G & A & T & C & T\\
			& & A & A & G & C & A & A & C & C\\
			& & T & T & G & G & A & A & C & T\\
			& & A & T & G & C & C & A & T & T\\
			& & A & T & G & G & C & A & C & T\\
			\hline
			\multirow{4}{*}{\textbf{Profilo}} & \textbf{A} & 5 & 1 & 0 & 0 & 5 & 5 & 0 & 0\\
			& \textbf{T} & 1 & 5 & 0 & 0 & 0 & 1 & 1 & 6\\
			& \textbf{G} & 1 & 1 & 6 & 3 & 0 & 1 & 0 & 0\\
			& \textbf{C} & 0 & 0 & 1 & 4 & 2 & 0 & 6 & 1\\
			\hline
			\textbf{Consenso} & & A & T & G & C & A & A & C & T 
		\end{tabular}
	\end{center}
	$Score(s,DNA) = 5+5+6+4+5+5+6+6=42$
\end{frame}

\begin{frame}{Score}
	Quanto può valere lo score?
	\begin{equation*}
		Score(s,DNA) = \begin{cases}
			l\cdot t, & \mbox{nel caso migliore}\\
			\frac{l\cdot t}{4}, & \mbox{nel caso peggiore}
		\end{cases}
	\end{equation*}
	\begin{itemize}
		\item $lt$ corrisponde al caso in cui tutti gli l-mer sono identici
		\item $\frac{lt}{4}$ corrisponde al caso in cui gli l-mer siano diversi in tutte le posizioni
	\end{itemize}
\end{frame}

\begin{frame}
  \frametitle{Algoritmi brute force}
\end{frame}

\begin{frame}
  \frametitle{Algoritmi greedy}
\end{frame}

\begin{frame}
  \frametitle{Algoritmi randomizzati}
\end{frame}

\begin{frame}
  \frametitle{Conclusioni}
\end{frame}

\begin{frame}
  \begin{center}
	\textbf{\calligra\Huge Fine.}\\
  \includegraphics[width=5cm]{img/ornament.eps}\\[1cm]
	\pause
	{\huge\calligra Domande?\pause{} Grazie!}
  \end{center}
\end{frame}
\end{document}